%! TEX root = main.tex

\begin{exercise}\label{ex:em-9}
Two long straight parallel wires are $8.0\ \mathrm{cm}$ apart. Wire~1 carries $I_1 = 6.0\ \mathrm{A}$, wire~2 carries $I_2 = 4.0\ \mathrm{A}$, in the same direction.
\begin{enumerate}[label=(\alph*)]
  \item Find the magnetic field produced by wire~1 at the position of wire~2 (magnitude and direction, described in words with axes you choose).
  \item Find the force per unit length on wire~2 (magnitude and direction). Is it attractive or repulsive?
  \item A solenoid has $N = 400$ turns over $\ell = 20\ \mathrm{cm}$ and carries $I = 3.0\ \mathrm{A}$. Find its interior field. How many times Earth's field ($\approx 50\ \mu\mathrm{T}$) is this?
\end{enumerate}
\end{exercise}

\begin{exercise}\label{ex:em-10}
A beam of singly ionized carbon isotopes ($q = +e$) with a spread of speeds enters a velocity selector with $E = 4.0\times 10^{4}\ \mathrm{V/m}$ and $B = 0.20\ \mathrm{T}$, then enters a spectrometer chamber with $B' = 0.20\ \mathrm{T}$ perpendicular to the beam. In the chamber each ion traces a half-circle and hits a detector plate; the distance from the entry slit to the impact point is the diameter $2r$.
\begin{enumerate}[label=(\alph*)]
  \item Show that only ions with one specific speed pass the selector, and compute that speed.
  \item Find the orbit radius in the spectrometer for $^{12}\mathrm{C}$ ($m = 12\,u$).
  \item The second isotope $^{13}\mathrm{C}$ lands $2.1\ \mathrm{cm}$ further out along the plate. Verify this separation from the formulas (no full recalculation needed: use scaling, and remember the plate measures diameters).
\end{enumerate}
\end{exercise}

\begin{exercise}\label{ex:em-11}
A power station delivers $P = 100\ \mathrm{MW}$ over a line of resistance $R = 2.0\ \Omega$, first at $U_\text{rms} = 20\ \mathrm{kV}$, then stepped up to $400\ \mathrm{kV}$.
\begin{enumerate}[label=(\alph*)]
  \item Find the rms current and the power lost as heat in the line in each case.
  \item Give the loss as a percentage of delivered power in each case.
  \item Explain in one or two sentences why the grid uses high voltage for transmission and steps it down near consumers, naming the device that makes this possible.
\end{enumerate}
\end{exercise}

\begin{exercise}\label{ex:em-12}
A laser beam ($\lambda = 532\ \mathrm{nm}$) strikes an air--glass boundary ($n_\text{glass} = 1.52$) at $50^\circ$ to the normal.
\begin{enumerate}[label=(\alph*)]
  \item Find the angle of the refracted ray inside the glass.
  \item The refracted ray hits the opposite face of a glass slab (glass to air). The slab faces are parallel. Show that the exiting ray is parallel to the incoming ray, and find the critical angle for the glass--air face. Does total internal reflection occur here?
  \item An optical fibre has core $n = 1.50$ and cladding $n = 1.48$. Find the critical angle at the core--cladding boundary and explain in words why light stays in the fibre.
\end{enumerate}
\end{exercise}

\begin{exercise}\label{ex:em-13}
A converging lens has focal length $f = 12\ \mathrm{cm}$.
\begin{enumerate}[label=(\alph*)]
  \item An object stands $a = 30\ \mathrm{cm}$ from the lens. Find the image distance $b$ and magnification $m$. Is the image real or virtual, upright or inverted?
  \item The object is moved to $a = 8.0\ \mathrm{cm}$ (inside the focal length). Find $b$ and $m$ and describe the image (this is the magnifying-glass case).
  \item A flat mirror stands $2.0\ \mathrm{m}$ from a person $1.7\ \mathrm{m}$ tall. How tall must the mirror be for the person to see their whole body, and at what height must its top edge sit? (Draw the rays; the answer is independent of the $2.0\ \mathrm{m}$.)
\end{enumerate}
\end{exercise}

\begin{exercise}\label{ex:em-14}
In a grating experiment the screen is $L = 1.20\ \mathrm{m}$ from a grating with $500$ lines per mm. A green laser produces first-order maxima $33.1\ \mathrm{cm}$ on either side of the central maximum.
\begin{enumerate}[label=(\alph*)]
  \item Find the slit spacing $d$ and the angle $\theta_1$ of the first maximum.
  \item Use $d\sin\theta_1 = \lambda$ to determine the laser wavelength.
  \item The laser is replaced by white light. Explain in words what pattern appears on the screen and why the centre stays white while higher orders smear into spectra.
\end{enumerate}
\end{exercise}

\begin{exercise}\label{ex:em-15}
An ideal transformer has $N_1 = 200$ primary turns and $N_2 = 20$ secondary turns. The primary is fed $U_{1,\text{rms}} = 230\ \mathrm{V}$ at $50\ \mathrm{Hz}$; the secondary drives a load $R = 4.0\ \Omega$.
\begin{enumerate}[label=(\alph*)]
  \item Starting from Faraday's law applied to each coil, derive $U_2/U_1 = N_2/N_1$ (state the shared-flux assumption explicitly) and find $U_{2,\text{rms}}$.
  \item Find the secondary current and the power delivered to the load. What primary current does energy conservation predict?
  \item Give two mechanisms that make a real transformer deliver slightly less than this, with one sentence each linking to the physics of this lesson.
\end{enumerate}
\end{exercise}

\begin{exercise}\label{ex:em-16}
A plane wave of wavelength $\lambda$ passes $N$ identical narrow slits spaced $d$ apart. The complex amplitude in direction $\theta$ is the phasor sum
\[
  A(\theta) = A_0\sum_{k=0}^{N-1} e^{ik\phi}, \qquad \phi = \frac{2\pi d\sin\theta}{\lambda},
\]
and the observed intensity is $I(\theta) = \abs{A(\theta)}^2$.
\begin{enumerate}[label=(\alph*)]
  \item Show that the principal maxima sit exactly at the grating condition $d\sin\theta = m\lambda$, for any $N$.
  \item Simulate $I(\theta)$ numerically (e.g.\ in Python) for $N = 2$ and $N = 8$ over $-90^\circ < \theta < 90^\circ$ with $d = 4\lambda$, and plot both. Count the small secondary maxima between neighbouring principal maxima for each $N$ and state the pattern.
  \item Explain in words why increasing $N$ makes the principal maxima sharper, connecting the phasor picture to the experimental fact that a grating resolves wavelengths far better than a double slit.
\end{enumerate}
\end{exercise}
